\section{Introduction}

\subsection{Background and Problem}

Accelerate is an embedded language in Haskell for high-performance array computations on GPUs and CPUs.
Its compiler translates high-level array operations into code for the target hardware, applying optimizations such as array fusion.
Fusion combines multiple array operations into a combined operation, reducing the need for intermediate arrays and thereby reducing memory allocations and data movement.

The compiler must decide which operations to fuse.
In the current ILP-based approach, these decisions are expressed as an integer linear programming (ILP) problem, which is solved using an external ILP solver.
The objective is to find the optimal fusion strategy that minimizes the cost, e.g. the number of memory allocations.
For some programs solving can take a long time and contributes significantly to the compilation time.

This project investigates how properties of the fusion problem can be used to reduce solving time.
The first step is to develop a domain-specific constraint language that makes these properties explicit.
This language then enables the implementation of a presolver that simplifies the problem before passing it to the external solver.
Building on the same representation, a custom solver was implemented that exploits the structure of the fusion problem as an alternative to the external solver.

\subsection{Scope and Contributions}

The project contributes a domain-specific constraint language for the fusion problem, together with a translation into an ILP representation for the external solver.
Building on this language, it introduces a presolver for simplifying fusion problems and a custom solver that exploits their structure.
The evaluation compares different presolving and solving configurations on a set of benchmarks.
CBC is used as an oracle to verify the correctness of the objective values returned by the custom solver.
The repositories containing the implementation and evaluation code are listed in Appendix~\ref{app:source-code}.

The report is structured as follows.
\begin{itemize}
    \item Section~\ref{sec:implementation} describes the implementation of the domain-specific constraint language, the presolver, and the custom solver.
    \item Section~\ref{sec:evaluation} describes the experimental evaluation with benchmarks, their results, and a discussion of them.
    \item Section~\ref{sec:conclusion} concludes the report.
\end{itemize}
